\chapter{Many-body theory with finitely many modes}

% B1 proposal (2026-11-26): concrete Fock space on configurations, bridged to the lattice code.
% Alternative considered: Physlib's WickAlgebra (abstract, QFT-oriented, no finite-mode Hilbert space);
% to be used later for Wick's theorem via its universality property, if our operators fit its hypotheses.

\begin{definition}[Fermionic Fock space of $M$ modes]\label{def:fock}
  \lean{QM.ManyBody.Fock, QM.ManyBody.FockOp, QM.ManyBody.finrank_fock}\leanok
  $\mathcal F_M$ is the Hilbert space of functions $\Psi\colon \mathcal P(\{0,\dots,M-1\})\to\mathbb C$ on all configurations
  (subsets of the modes), with $\langle\Psi,\Phi\rangle=\sum_s\overline{\Psi(s)}\Phi(s)$; $\dim\mathcal F_M = 2^M$.
  The basis vector $|s\rangle$ is the indicator of the configuration $s$.
\end{definition}

\begin{definition}[Creation and annihilation operators]\label{def:cdag}
  \lean{QM.ManyBody.jwSign, QM.ManyBody.cre, QM.ManyBody.ann, QM.ManyBody.ann_eq_adjoint,
    QM.ManyBody.cre_mulVec_single, QM.ManyBody.ann_mulVec_single,
    QM.ManyBody.creOp, QM.ManyBody.annOp, QM.ManyBody.annOp_eq_adjoint}\leanok
  \uses{def:fock}
  With the Jordan–Wigner sign $\sigma_i(s)=(-1)^{|\{j\in s:\ j<i\}|}$:
  $c_i^\dagger|s\rangle=\sigma_i(s)\,|s\cup\{i\}\rangle$ if $i\notin s$ and $0$ otherwise;
  $c_i=(c_i^\dagger)^\dagger$, so $c_i|s\rangle=\sigma_i(s)\,|s\setminus\{i\}\rangle$ if $i\in s$ and $0$ otherwise.
\end{definition}

\begin{theorem}[Canonical anticommutation relations]\label{thm:car}
  \lean{QM.ManyBody.ann_mul_cre_add, QM.ManyBody.ann_mul_ann_add, QM.ManyBody.cre_mul_cre_add,
    QM.ManyBody.annOp_mul_creOp_add, QM.ManyBody.annOp_mul_annOp_add, QM.ManyBody.creOp_mul_creOp_add}\leanok
  \uses{def:cdag}
  For all modes $i,j$: $\{c_i,c_j^\dagger\}=\delta_{ij}\,\mathbb 1$ and $\{c_i,c_j\}=\{c_i^\dagger,c_j^\dagger\}=0$.
\end{theorem}
\begin{proof}\leanok
  Compare the action on vectors. Adding an unoccupied mode $j$ multiplies $\sigma_i$ by $-1$ exactly when $j<i$; for
  $i\neq j$ the two exchange factors multiply to $-1$, so the two terms of each anticommutator cancel. For $i=j$,
  $\sigma_i(s)^2=1$ and exactly one of the two terms of $\{c_i,c_i^\dagger\}$ survives on each configuration.
  $\{c_i^\dagger,c_j^\dagger\}=0$ is the conjugate transpose of $\{c_j,c_i\}=0$. The identities hold first for the
  matrices indexed by configurations and are transported to bounded operators on $\mathcal F_M$ by the star-algebra
  equivalence between matrices and operators (\texttt{Matrix.toEuclideanCLM}).
\end{proof}

\begin{theorem}[Number operator]\label{thm:number}
  \lean{QM.ManyBody.num, QM.ManyBody.numTot, QM.ManyBody.num_mulVec_single, QM.ManyBody.num_mulVec_single_conf,
    QM.ManyBody.numTot_mulVec_single, QM.ManyBody.num_isHermitian, QM.ManyBody.numTot_isHermitian,
    QM.ManyBody.num_mul_num, QM.ManyBody.numTotOp_isSelfAdjoint}\leanok
  \uses{def:cdag}
  $\hat N=\sum_i c_i^\dagger c_i$ acts as $\hat N|s\rangle=|s|\,|s\rangle$; $\hat n_i=c_i^\dagger c_i$ acts as the occupation
  $\mathtt{occF}\ i\ s$ of the lattice code (\texttt{QM.ManyBody.Configurations}).
\end{theorem}
\begin{proof}\leanok
  $c_i^\dagger c_i|s\rangle$: $c_i$ removes $i$ with the sign $\sigma_i(s\setminus\{i\})$ and $c_i^\dagger$ puts it back with the
  same sign, whose square is $1$; if $i\notin s$ the result is $0$. Summing the indicators $[i\in s]$ over $i$ gives $|s|$.
  Both operators are self-adjoint, and $\hat n_i^2=\hat n_i$.
\end{proof}

% Particle-hole statements written 2026-11-26 while finishing B1, to be reviewed by Miguel.
\begin{theorem}[Particle–hole transformation]\label{thm:particle-hole}
  \lean{QM.ManyBody.ph, QM.ManyBody.ph_mulVec_single, QM.ManyBody.card_compl_conf, QM.ManyBody.ph_mul_ph,
    QM.ManyBody.ph_isHermitian, QM.ManyBody.ph_mul_conjTranspose, QM.ManyBody.jwSign_compl,
    QM.ManyBody.ph_cre_ph, QM.ManyBody.ph_ann_ph, QM.ManyBody.ph_num_ph, QM.ManyBody.ph_numTot_ph,
    QM.ManyBody.stag, QM.ManyBody.ph_oneBody_ph, QM.ManyBody.phOp, QM.ManyBody.phOp_isSelfAdjoint,
    QM.ManyBody.phOp_mul_phOp, QM.ManyBody.phOp_creOp_phOp}\leanok
  \uses{def:cdag, thm:car, thm:number}
  Let $P|s\rangle=|s^c\rangle$. Then $P$ is a self-adjoint unitary with $P^2=1$, and it maps the $N$-particle sector onto
  the $(M-N)$-particle sector. It acts on the operators as
  $Pc_i^\dagger P=(-1)^i c_i$, $Pc_iP=(-1)^i c_i^\dagger$, $P\hat n_iP=1-\hat n_i$ and $P\hat NP=M-\hat N$. For a one-body
  operator, $P\big(\sum_{ij}h_{ij}c_i^\dagger c_j\big)P=(\operatorname{tr}h)\,\mathbb 1-\sum_{ij}\tilde h_{ji}\,c_i^\dagger c_j$
  with the staggered matrix $\tilde h_{ij}=(-1)^{i+j}h_{ij}$.
\end{theorem}
\begin{proof}\leanok
  Every mode below $i$ lies either in $s$ or in $s^c$, so $\sigma_i(s)\,\sigma_i(s^c)=(-1)^i$. Comparing the action on vectors
  gives $Pc_i^\dagger=(-1)^ic_iP$, and $P^2=1$ turns this into $Pc_i^\dagger P=(-1)^ic_i$. The conjugate transpose gives the
  relation for $c_i$. Then $Pc_i^\dagger c_jP=(-1)^{i+j}c_ic_j^\dagger=(-1)^{i+j}(\delta_{ij}-c_j^\dagger c_i)$ by the
  anticommutation relations, which gives the number operators and the one-body operators.
\end{proof}

\begin{theorem}[Bridge to the lattice code]\label{thm:fock-bridge}
  \lean{QM.ManyBody.oneBody, QM.ManyBody.oneBody_apply_eq_zero_of_card_ne, QM.ManyBody.slaterVec,
    QM.ManyBody.hop_mulVec_slaterVec, DFT.Lattice.Fermions.slaterVec_apply, DFT.Lattice.Fermions.slaterVec_τs,
    DFT.Lattice.Fermions.oneBody_apply_conf, DFT.Lattice.Fermions.fockHam_apply_conf,
    DFT.Lattice.Fermions.H0F_eq_fock, DFT.Lattice.Fermions.hamN_H0F_eq_fock}\leanok
  \uses{def:cdag, thm:car}
  For a real single-particle matrix $h$, the operator $\sum_{ij}h_{ij}\,c_i^\dagger c_j$ preserves the $N$-particle sector
  $\{s:|s|=N\}$ and its matrix there is the many-body matrix $\Gamma\,h$ of \texttt{DFT/Lattice/Fermions/SecondQuantization}
  (\texttt{Conf M N} ordered as in that file). Consequently the lattice Hamiltonians \texttt{H0F} are restrictions of
  second-quantized Hamiltonians on $\mathcal F_M$.
\end{theorem}
\begin{proof}\leanok
  From the anticommutation relations, $c_i^\dagger c_j c_a^\dagger=\delta_{ja}c_i^\dagger+c_a^\dagger c_i^\dagger c_j$. By
  induction on $N$, $c_i^\dagger c_j$ acts on a creation string $c_{\tau_0}^\dagger\cdots c_{\tau_{N-1}}^\dagger|\emptyset\rangle$
  as a derivation. It replaces each occurrence of the mode $j$ in $\tau$ by $i$. Expanding the determinant $\mathtt{amp}\,s\,\tau$
  along its first row shows, again by induction on $N$, that the coefficient of the creation string on $s$ is
  $\mathtt{amp}\,s\,\tau$. The sign $(-1)^k$ of the cofactor is the Jordan–Wigner sign, because the $k$-th occupied mode of
  $s$ has exactly $k$ occupied modes below it. For the sorted modes of $s'$ the creation string is $|s'\rangle$. Hence
  $\langle s|\sum_{ij}h_{ij}c_i^\dagger c_j|s'\rangle=\sum_k\sum_i h_{i,\tau_k}\,\mathtt{amp}\,s\,(\tau\text{ with }\tau_k\to i)$,
  which is $(\Gamma\,h^{T})_{s's}=(\Gamma\,h)_{ss'}$. Each term $c_i^\dagger c_j$ removes one particle and adds one, so the
  matrix elements between different particle numbers vanish.
\end{proof}

\begin{corollary}[The lattice one-body reduced density matrix]\label{cor:fock-rdm}
  \lean{DFT.rdm, DFT.Lattice.RDMFT.toFock, DFT.Lattice.RDMFT.oneBody_single, DFT.Lattice.RDMFT.rdm_eq_inner}\leanok
  \uses{thm:fock-bridge}
  Let $\psi$ be a real $N$-particle state on the lattice and $\Psi\in\mathcal F_M$ its extension by zero outside the
  $N$-particle sector. The one-body reduced density matrix of the lattice code, $\gamma_{ij}=\langle\psi,\Gamma(E_{ij})\psi\rangle$,
  is $\gamma_{ij}=\langle\Psi,c_i^\dagger c_j\Psi\rangle$.
\end{corollary}
\begin{proof}\leanok
  $\sum_{kl}(E_{ij})_{kl}c_k^\dagger c_l=c_i^\dagger c_j$. Both sums over configurations reduce to the $N$-particle sector,
  where Theorem~\ref{thm:fock-bridge} identifies the matrix elements of $c_i^\dagger c_j$ with those of $\Gamma(E_{ij})$.
\end{proof}

% B2 statements written and approved by Miguel 2026-11-26. Convention chosen by Miguel:
% gamma_ij = <c_i^dag c_j>, as in the lattice code (Corollary cor:fock-rdm). Coleman's theorem belongs to B2.

\begin{definition}[Reduced density matrices]\label{def:rdm12}
  \lean{QM.ManyBody.rdm1, QM.ManyBody.rdm2, QM.ManyBody.IsNState, QM.ManyBody.IsNEnsemble,
    QM.ManyBody.rdm1Ens, QM.ManyBody.rdm2Ens}\leanok
  \uses{def:cdag}
  For $\Psi\in\mathcal F_M$ the one-body reduced density matrix is $\gamma_{ij}(\Psi)=\langle\Psi,c_i^\dagger c_j\Psi\rangle$ and the
  two-body reduced density matrix is $D_{ij,kl}(\Psi)=\langle\Psi,c_i^\dagger c_j^\dagger c_l c_k\Psi\rangle$. The density is
  $\rho_i(\Psi)=\gamma_{ii}(\Psi)$. An $N$-particle state is a unit vector with $\hat N\Psi=N\Psi$. An $N$-particle ensemble is a
  finite family of $N$-particle states $\Psi_k$ with weights $w_k\ge0$, $\sum_kw_k=1$. Its reduced density matrices are the
  weighted sums $\sum_kw_k\gamma(\Psi_k)$ and $\sum_kw_kD(\Psi_k)$.
\end{definition}

\begin{theorem}[Necessary conditions on the one-body matrix]\label{thm:rdm1-necessary}
  \lean{QM.ManyBody.quad_rdm1, QM.ManyBody.quad_one_sub_rdm1, QM.ManyBody.rdm1_isHermitian,
    QM.ManyBody.rdm1_posSemidef, QM.ManyBody.one_sub_rdm1_posSemidef, QM.ManyBody.trace_rdm1,
    QM.ManyBody.rdm1Ens_isHermitian, QM.ManyBody.rdm1Ens_posSemidef, QM.ManyBody.one_sub_rdm1Ens_posSemidef,
    QM.ManyBody.trace_rdm1Ens}\leanok
  \uses{def:rdm12, thm:car, thm:number}
  For every $N$-particle state or ensemble, $\gamma$ is Hermitian, $0\le\gamma\le1$ as a matrix (both $\gamma$ and $1-\gamma$ are
  positive semidefinite) and $\operatorname{tr}\gamma=N$.
\end{theorem}
\begin{proof}\leanok
  With $c(x)=\sum_jx_jc_j$ one has $\sum_{ij}\overline{x_i}x_j\gamma_{ij}=\|c(x)\Psi\|^2\ge0$. By the anticommutation relations
  $c(x)c(x)^\dagger+c(x)^\dagger c(x)=\|x\|^2$, so $\|x\|^2-\sum_{ij}\overline{x_i}x_j\gamma_{ij}=\|c(x)^\dagger\Psi\|^2\ge0$. The trace
  is $\langle\Psi,\hat N\Psi\rangle=N$. All three properties are preserved by convex combinations.
\end{proof}

\begin{corollary}[N-representability of the density]\label{cor:density-rep}
  \lean{DFT.Lattice.RDMFT.density_nrep, DFT.Lattice.RDMFT.density_nrep_nonneg, DFT.Lattice.RDMFT.rdm1_toFock,
    DFT.Lattice.RDMFT.rdm_diag, DFT.Lattice.RDMFT.toFock_isNState}\leanok
  \uses{thm:rdm1-necessary, cor:fock-rdm}
  For $\rho\colon\{0,\dots,M-1\}\to\mathbb R$ the following are equivalent.
  (a) $\rho$ is the density of an $N$-particle state.
  (b) $\rho$ is the density of an $N$-particle ensemble.
  (c) $0\le\rho_i\le1$ for every $i$ and $\sum_i\rho_i=N$.
  The state in (a) can be chosen with nonnegative real coefficients on configurations.
\end{corollary}
\begin{proof}\leanok
  (a)$\Rightarrow$(b) is trivial and (b)$\Rightarrow$(c) is the diagonal of Theorem~\ref{thm:rdm1-necessary}. For
  (c)$\Rightarrow$(a), $\rho$ is a mixture $\sum_s\mu_s\,\mathrm{occ}(s)$ of configuration densities by the vertex description
  of the hypersimplex (\texttt{DFT.hypersimplex\_mixture}). The state $\Psi=\sum_s\sqrt{\mu_s}\,|s\rangle$ has density $\rho$,
  because $c_i^\dagger c_i$ is diagonal on configurations.
\end{proof}

\begin{theorem}[Basic properties of the two-body matrix]\label{thm:rdm2-basic}
  \lean{QM.ManyBody.num_mul_ann, QM.ManyBody.numTot_mul_ann, QM.ManyBody.numTotOp_annOp_of,
    QM.ManyBody.rdm2_swap_left, QM.ManyBody.rdm2_swap_right, QM.ManyBody.rdm2_isHermitian,
    QM.ManyBody.sum_rdm2_contract, QM.ManyBody.sum_rdm2_trace, QM.ManyBody.rdm2Ens_isHermitian,
    QM.ManyBody.rdm2Ens_swap, QM.ManyBody.sum_rdm2Ens_contract, QM.ManyBody.sum_rdm2Ens_trace}\leanok
  \uses{def:rdm12, thm:car, thm:number}
  For every $N$-particle state or ensemble,
  $D_{ji,kl}=-D_{ij,kl}=D_{ij,lk}$ and $D_{kl,ij}=\overline{D_{ij,kl}}$. The contraction is
  $\sum_kD_{ik,jk}=(N-1)\,\gamma_{ij}$ and the trace is $\sum_{ij}D_{ij,ij}=N(N-1)$.
\end{theorem}
\begin{proof}\leanok
  Antisymmetry follows from $\{c_i^\dagger,c_j^\dagger\}=0$ and $\{c_k,c_l\}=0$, Hermiticity from taking adjoints. For the
  contraction, $\sum_kc_k^\dagger c_kc_j=\hat Nc_j=c_j(\hat N-1)$, and $c_j$ lowers the particle number by one. The trace
  follows by contracting once more.
\end{proof}

\begin{theorem}[The P, Q and G conditions]\label{thm:pqg}
  \lean{QM.ManyBody.quad_rdm2, QM.ManyBody.rdm2_posSemidef, QM.ManyBody.qMat, QM.ManyBody.qMat_eq_inner,
    QM.ManyBody.quad_qMat, QM.ManyBody.qMat_posSemidef, QM.ManyBody.gMat, QM.ManyBody.gMat_eq_inner,
    QM.ManyBody.quad_gMat, QM.ManyBody.gMat_posSemidef, QM.ManyBody.rdm2Ens_posSemidef,
    QM.ManyBody.qMatEns_posSemidef, QM.ManyBody.gMatEns_posSemidef}\leanok
  \uses{thm:rdm2-basic}
  For every $N$-particle state or ensemble, the three Hermitian matrices indexed by pairs of modes
  \begin{align*}
    P_{ij,kl}&=D_{ij,kl},\\
    Q_{ij,kl}&=\delta_{ik}\delta_{jl}-\delta_{il}\delta_{jk}-\delta_{jl}\gamma_{ki}+\delta_{jk}\gamma_{li}+\delta_{il}\gamma_{kj}
      -\delta_{ik}\gamma_{lj}+D_{kl,ij},\\
    G_{ij,kl}&=\delta_{jl}\gamma_{ik}-D_{il,kj}
  \end{align*}
  are positive semidefinite.
\end{theorem}
\begin{proof}\leanok
  $P$, $Q$ and $G$ are the expectation values $\langle c_i^\dagger c_j^\dagger c_lc_k\rangle$, $\langle c_ic_jc_l^\dagger c_k^\dagger\rangle$
  and $\langle c_i^\dagger c_jc_l^\dagger c_k\rangle$. Their quadratic forms are $\|A\Psi\|^2$ with $A=\sum x_{kl}c_lc_k$,
  $\sum x_{kl}c_l^\dagger c_k^\dagger$ and $\sum x_{kl}c_l^\dagger c_k$ respectively. Normal ordering with the anticommutation
  relations expresses $Q$ and $G$ through $D$ and $\gamma$ as stated.
\end{proof}

\begin{definition}[Mode rotations]\label{def:mode-rotation}
  \uses{def:cdag}
  For a unitary $u\in U(M)$ there is a unitary $U(u)$ on $\mathcal F_M$ with $U(u)\,c_i^\dagger\,U(u)^\dagger=\sum_ju_{ji}c_j^\dagger$
  for every $i$ and $U(u)|\emptyset\rangle=|\emptyset\rangle$. It commutes with $\hat N$, $U(uv)=U(u)U(v)$, and
  $\gamma(U(u)\Psi)=\bar u\,\gamma(\Psi)\,u^{T}$.
\end{definition}

\begin{theorem}[Coleman]\label{thm:coleman}
  \uses{thm:rdm1-necessary, def:mode-rotation}
  A complex $M\times M$ matrix $\gamma$ is the one-body reduced density matrix of an $N$-particle ensemble if and only if
  $\gamma$ is Hermitian, $0\le\gamma\le1$ and $\operatorname{tr}\gamma=N$.
\end{theorem}
\begin{proof}
  Necessity is Theorem~\ref{thm:rdm1-necessary}. Conversely, diagonalize $\gamma=\bar u\,\operatorname{diag}(\lambda)\,u^{T}$
  with $u$ unitary. The eigenvalues satisfy $0\le\lambda_i\le1$ and $\sum_i\lambda_i=N$, so
  $\lambda=\sum_s\mu_s\,\mathrm{occ}(s)$ by the hypersimplex. The ensemble of configurations $|s\rangle$ with weights $\mu_s$ has
  $\gamma=\operatorname{diag}(\lambda)$. Applying $U(u)$ to each state gives $\gamma$.
\end{proof}

\begin{definition}[Gibbs state]\label{def:gibbs}
  $\rho_\beta = e^{-\beta(H-\mu N)}/Z$ on $\mathcal F_M$.
\end{definition}

\begin{theorem}[KMS condition]\label{thm:kms}
  \uses{def:gibbs}
  $\operatorname{Tr}\rho_\beta\, A(t) B = \operatorname{Tr}\rho_\beta\, B A(t+i\beta)$ (finite dimensions: an identity of
  entire functions of $t$).
\end{theorem}
